\documentclass[11pt]{article}
\usepackage[margin=1in]{geometry}
\usepackage[T1]{fontenc}
\usepackage{lmodern}
\usepackage{microtype}
\usepackage{setspace}
\usepackage{amsmath}
\usepackage{amssymb}
\usepackage{booktabs}
\usepackage{longtable}
\usepackage{array}
\usepackage{hyperref}
\hypersetup{colorlinks=true,linkcolor=blue,urlcolor=blue,citecolor=blue}
\setstretch{1.12}

\title{Sparse Signal-Field SAE Project: Research Report on Planned vs Executed Experiments}
\author{Technical Project Report}
\date{June 1, 2026}

\begin{document}
\maketitle

\begin{abstract}
This report explains what the project planned to run, what was actually implemented and executed, and what the resulting evidence supports. The original design was an eight-cell matrix that crossed training protocol, SAE architecture, and mask type. Execution completed the full Family I matrix (E1--E4), then focused on winner-path hierarchy experiments and chain stabilization. Family II was therefore executed partially in modified form, not as full E5--E8 matrix coverage. Quantitative results show strong sparse reconstruction sufficiency and strong causal upper-bound performance, while raw multi-step learned chains collapse unless each step is re-grounded on the SAE manifold. This document is written as a research report with explicit terminology, method statements, execution accounting, and result interpretation.
\end{abstract}

\section{Objective and Reporting Scope}
The reporting objective is to give a technically precise account of the current project state without requiring the reader to infer missing context from code files. The report addresses four concrete questions. The first question is experiment coverage: which of E1--E8 were actually executed. The second question is implementation drift: what changed between the high-level design and the real code path. The third question is execution summary: what experiments were run in practice after those changes. The fourth question is numerical evidence: what the recorded outputs show.

All statements are tied to artifacts in \texttt{Vision\_Sparsity/CNN-SAE}, including the high-level design document, source files, run directories, and result JSON logs. The key planning artifact is \texttt{high-level/sparse\_signal\_sae\_experiment\_plan\_high\_accuracy\_conv.tex} with matching PDF. The main implementation artifacts are in \texttt{src/} and \texttt{scripts/}. The main execution artifacts are in \texttt{runs/} and \texttt{logs/}.

\section{Terminology and Experimental Matrix}
The project uses five tapped sections of a frozen CNN, denoted Q1 through Q5. At each section, an SAE maps activation field $H \in \mathbb{R}^{C\times H\times W}$ to sparse coefficient field $Z \in \mathbb{R}^{K\times H\times W}$, masking is applied in coefficient space, and a decoder reconstructs $\hat{H}$ for downstream evaluation through the frozen model tail.

The planned matrix in the high-level document is defined by three binary factors. The first factor is training protocol. Family I means independent layerwise SAE training followed by learned transitions between frozen code spaces. Family II means incremental dependent sparse-code learning. The second factor is SAE type: Field SAE versus Vector SAE. The third factor is mask type: global TopK versus receptive-field-local TopK.

This yields eight planned experiments:
\begin{center}
\begin{tabular}{llll}
\toprule
ID & Protocol & SAE type & Mask type \\
\midrule
E1 & Family I & Field & Global \\
E2 & Family I & Field & RF-local \\
E3 & Family I & Vector & Global \\
E4 & Family I & Vector & RF-local \\
E5 & Family II & Field & Global \\
E6 & Family II & Field & RF-local \\
E7 & Family II & Vector & Global \\
E8 & Family II & Vector & RF-local \\
\bottomrule
\end{tabular}
\end{center}

\section{Implemented System and Method Changes}
The executed codebase implemented both SAE families and both mask families, then developed an additional chain-stability mechanism that became central to later results. In \texttt{src/sae.py}, Field SAE is implemented as a spatially expressive local-residual convolutional encoder/decoder, while Vector SAE is implemented as shared per-location 1x1 mappings with decoder normalization. In \texttt{src/masks.py}, global TopK and RF-local TopK are both implemented with hard masking.

The training workflow in \texttt{src/train\_sae.py} uses a staged curriculum (warmup, sparsity anneal, fixed budget), which is a concrete implementation refinement relative to a naive fixed-sparsity start. Transition training is implemented in \texttt{src/train\_transition.py} and \texttt{src/transitions.py}. The same transition module also implements the frozen-block hybrid operator, which uses real frozen backbone blocks between decoded and re-encoded sparse states to measure a causal upper bound independent of learned transition error.

The most important implementation change relative to the original matrix intent appears in chain handling. \texttt{src/eval\_chain.py} evaluates raw chaining, remasked chaining, re-grounded chaining, and hybrid chaining. \texttt{src/train\_chain.py} implements chain-aware training with scheduled sampling and optional BPTT through re-grounding. This is a real protocol extension, but it is applied on the selected winner configuration rather than as full E5--E8 matrix breadth.

\section{Execution Accounting: What Was Actually Run}
Execution records show full run coverage for E1 through E4 and no directory-level matrix completion for E5 through E8. The artifact pattern is explicit in \texttt{runs/phase1} and \texttt{runs/phase2}: E3 appears in phase1, and E1/E2/E4 appear in phase2 across Q1--Q5. Corresponding result files exist for those cells.

For Family II, no full E5--E8 matrix directories exist. However, phase3b artifacts exist and are substantive. Specifically, \texttt{runs/phase3b/result.json} and \texttt{runs/phase3b\_bptt/result.json} show chain-aware training runs on frozen winner SAEs with scheduled sampling and re-grounding. Therefore, Family II-style mechanism work was executed, but matrix breadth was not completed.

The correct accounting statement is therefore two-dimensional. Breadth across planned Family II configuration cells is incomplete (1 of 4 configurations executed in Family II-style form, approximately 25\% breadth coverage). Depth along that one executed winner configuration is substantial because the full Q1$\rightarrow$Q5 chain is trained and evaluated.

\section{Chronological Execution Summary}
Phase 1 established the anchor condition using E3 (Vector, global, independent). This phase answered the feasibility question and included controls such as random-dictionary SAE runs and PCA baselines. Phase 2 branched to E1, E2, and E4 to compare architecture and masking effects under matched section coverage. Phase 3 trained Family I adjacent transitions on winner SAEs and evaluated chain behavior under multiple rollout modes. Phase 3b introduced chain-aware training with scheduled sampling and re-grounding, then extended this with BPTT through re-grounding.

Additional tier experiments were executed for intervention and feature-quality diagnostics. \texttt{runs/intervene\_result.json} and \texttt{runs/audit/audit\_Q3.json}, \texttt{runs/audit/audit\_Q5.json} are present and complete. Additional tier campaigns for taxonomy, ViT transfer, R20/C10 grid, seed grid, and R110 reconstruction were attempted but blocked by runtime CUDA/node state issues documented in \texttt{runs/tier\_rest} logs and reflected by zero successful outputs in the corresponding grid summaries.

\section{Results}
\subsection{Backbone and Reconstruction Baselines}
The frozen ResNet-56/CIFAR-100 reference is 71.83\% top-1. This value appears in \texttt{logs/backbone\_r56\_c100\_\_done.json}. Under sparse reconstruction conditions, E1 and E3 comparison at 2\% retained shows that Field+Global is materially stronger than Vector+Global in early sections, while differences narrow in deeper sections.

\begin{center}
\begin{tabular}{lcc}
\toprule
Section & E1 (Field+Global, 2\%) & E3 (Vector+Global, 2\%) \\
\midrule
Q1 & 0.7139 & 0.6376 \\
Q2 & 0.7164 & 0.6327 \\
Q3 & 0.7078 & 0.6646 \\
Q4 & 0.7106 & 0.7098 \\
Q5 & 0.7147 & 0.7123 \\
\bottomrule
\end{tabular}
\end{center}

RF-local comparisons at \texttt{k\_rf=4} show Field+RF remains viable while Vector+RF degrades sharply in several sections.

\begin{center}
\begin{tabular}{lcc}
\toprule
Section & E2 (Field+RF, k=4) & E4 (Vector+RF, k=4) \\
\midrule
Q1 & 0.7158 & 0.6858 \\
Q2 & 0.6992 & 0.4645 \\
Q3 & 0.6698 & 0.2430 \\
Q4 & 0.6949 & 0.5942 \\
Q5 & 0.6443 & 0.2504 \\
\bottomrule
\end{tabular}
\end{center}

\subsection{Transition and Chain Results}
Family I adjacent transitions in \texttt{runs/phase3/T\_*} show strong one-step performance on easier edges and a pronounced Q3$\rightarrow$Q4 gap.

\begin{center}
\begin{tabular}{lcc}
\toprule
Transition & Learned top-1 & Hybrid top-1 \\
\midrule
Q1$\rightarrow$Q2 & 0.7146 & 0.7151 \\
Q2$\rightarrow$Q3 & 0.6907 & 0.7105 \\
Q3$\rightarrow$Q4 & 0.6504 & 0.7113 \\
Q4$\rightarrow$Q5 & 0.7089 & 0.7149 \\
\bottomrule
\end{tabular}
\end{center}

Chain rollout in \texttt{runs/phase3/chain\_result.json} separates information sufficiency from learned-chain stability. Raw learned chain is 0.0116, remasked chain is 0.2126, re-grounded chain is 0.5930, and hybrid chain is 0.7122 versus original 0.7183. This indicates that sparse support is information-sufficient under real frozen propagation but unstable under naive recursive learned-code propagation.

\subsection{Phase 3b Family II-style Results}
Phase3b scheduled-sampling run gives re-grounded chain top-1 of 0.6624 with hybrid 0.7122. Phase3b BPTT variant gives re-grounded chain top-1 of 0.7010 with hybrid still 0.7122. Raw chain remains near chance in both phase3b variants. These outcomes support the interpretation that manifold re-grounding is the key mechanism for stable multi-step sparse-code composition.

\subsection{Causal and Feature-Quality Diagnostics}
Intervention results show top-vs-random ablation impact ratios above one with large margins on most edges, including approximately 3.29x and 7.5x in reported pairs, with locality concentration above null area baselines. Feature audits report low dead-feature and duplicate rates at Q3 and Q5. These diagnostics support that the sparse representation is structured and nontrivial rather than degenerate.

\section{Family I vs Family II: Correct Final Interpretation}
The final interpretation must separate matrix completion from mechanism execution. E1 through E4 were fully executed as planned matrix cells. E5 through E8 were not executed as full matrix cells. At the same time, Family II-style training was executed in modified form on the winner configuration through phase3b and phase3b\_bptt. Therefore, the project did run Family II-style methods, but did not complete planned Family II breadth.

This interpretation avoids two opposite errors. The first error is claiming ``Family II was not run,'' which conflicts with phase3b artifacts. The second error is claiming ``Family II was fully run,'' which conflicts with missing E6--E8 breadth and absence of full E5--E8 matrix outputs.

\section{Limitations and Pending Work}
The main scientific limitation is incomplete Family II breadth. The main operational limitation is blocked tier execution under CUDA/node instability. Grid summary files for \texttt{runs/r20c10} and \texttt{runs/seeds} show 0 successful result outputs under attempted runs. Taxonomy and ViT transfer scripts are implemented but corresponding completion artifacts are absent under the blocked node state.

These limitations constrain claim scope. Strong claims are justified for executed winner-path mechanisms and Family I matrix outcomes. Broad claims over all planned Family II architecture/mask combinations are not yet justified.

\section{Conclusion}
This project established strong evidence that frozen CNN activations are sparse-compressible in a learned transform domain with high downstream fidelity at low retained fractions. It further established that naive learned sparse-code chains are unstable under recursive rollout, while re-grounding and chain-aware training recover most of the gap to a causal upper bound defined by frozen-block hybrid propagation.

The implementation record is therefore coherent but not matrix-complete: Family I matrix coverage is complete, Family II-style mechanism coverage is substantial on one winner path, and full Family II matrix breadth remains future work. This is the accurate research-state summary supported by the current artifact set.

\section*{Artifact Trace}
Primary files for verification are \texttt{high-level/sparse\_signal\_sae\_experiment\_plan\_high\_accuracy\_conv.tex}, \texttt{scripts/run\_grid.py}, \texttt{src/train\_transition.py}, \texttt{src/eval\_chain.py}, \texttt{src/train\_chain.py}, \texttt{runs/phase1}, \texttt{runs/phase2}, \texttt{runs/phase3}, \texttt{runs/phase3b}, \texttt{runs/phase3b\_bptt}, \texttt{runs/intervene\_result.json}, \texttt{runs/audit/audit\_Q3.json}, \texttt{runs/audit/audit\_Q5.json}, and \texttt{logs/baselines\_pca.json}.

\end{document}
